\originalsection{作业 11: 角动量}
\begin{exercise}
原题 1. 半径 $R$、质量 $m_1$ 的环平放光滑桌面, 环缘一点固定为枢轴. 质量 $m_2$ 的虫从枢轴开始沿环逆时针以相对环速率 $v$ 走, 初始环静止, $\uv k$ 出页. (a) 虫走半圈时环角速度. (b) 虫回枢轴时环角速度.
\end{exercise}
\begin{solution}
枢轴角动量恒为零. (a) 半圈时虫距枢轴 $2R$, 相对环的切向速度方向使正角动量为 $2m_2Rv$. 环与虫的共同转动惯量为 $2m_1R^2+4m_2R^2$. 所以 $\vct\omega=-m_2v\uv k/[R(m_1+2m_2)]$. (b) 虫在枢轴处力臂零, 其角动量为零, 环也须回到 $\vct\omega=0$. 虫用内能做功, 能量不守恒.
\end{solution}
\begin{exercise}
原题 2. 均匀杆长 $d$、质量 $m$, 一端 P 固定, 静止于光滑水平桌面. 两个点质量各为 $m$, 一以 $v$ 向右撞距 P 为 $2d/3$ 处, 另一以 $v$ 向左同时撞距 P 为 $d/3$ 处, 均粘住. 两碰点在 P 同侧. (a) 取顺时针为正求末角速度. (b) 求 $(K_f-K_i)/K_i$.
\end{exercise}
\begin{solution}
初始角动量为 $mv(2d/3-d/3)=mvd/3$. 总惯量为 $md^2/3+m(2d/3)^2+m(d/3)^2=8md^2/9$. (a) $\omega_f=3v/(8d)$. (b) $K_i=mv^2$, $K_f=\frac12I\omega_f^2=mv^2/16$, 所以比为 $-15/16$. 固定枢轴有外冲量, 总线动量不守恒.
\end{solution}
\begin{exercise}
原题 3. 同质量 $m$ 的点物以 $v_i$ 向右弹性撞均匀杆中点, 杆长 $d$、一端固定, 桌面光滑. 物体向左反弹, 杆顺时针转动. 求杆末角速度.
\end{exercise}
\begin{solution}
令反弹速率 $v_f>0$, 枢轴角动量给 $mv_i d/2=-mv_f d/2+(md^2/3)\omega_f$. 能量为 $mv_i^2/2=mv_f^2/2+(md^2/6)\omega_f^2$. 非零冲量解 $v_f=v_i/7$, $\omega_f=12v_i/(7d)$. 也可把杆接触点等效质量取 $I/(d/2)^2=4m/3$, 但必须来自同样的支点约束.
\end{solution}
\begin{exercise}
原题 4. 质量 $m_R$、半径 $R$ 环平放光滑桌面, 环最上缘 P 固定. 点质量 $m$ 从左以 $v_0$ 向右撞环最左侧中点, 弹性碰撞后物体向左反弹、环逆时针. 求反弹速率.
\end{exercise}
\begin{solution}
力作用线到 P 距离 $R$, 环惯量 $I_P=2m_RR^2$. 令左反弹速率 $v_f$, 则 $mR(v_0+v_f)=I_P\omega_f$. 配合动能守恒解得 $v_f=(2m_R-m)v_0/(2m_R+m)$. 原题反弹方向要求 $2m_R>m$; 等号时物体停止, 小于时向右继续而不应把负值称为反弹速率.
\end{solution}
\begin{exercise}
原题 5. 航天器 1 在质量 $m_p$ 行星的半径 $R$ 圆轨道, 航天器 2 处于近日 $R$、远日 $3R$ 椭圆轨道. 行星质量远大于两者. 两者同平面同方向, 到航天器 2 的近日点时它点火调整速度以与 1 对接, 时间很短. 求 $\Delta v=v_1-v_p$, 只含 $G,R,m_p$, 并判断加速或减速.
\end{exercise}
\begin{solution}
椭圆半长轴 $a=2R$, 能量 $E/m=-Gm_p/(4R)$. 近日能量给 $v_p^2=2Gm_p/R-Gm_p/(2R)=3Gm_p/(2R)$. 圆速 $v_1=\sqrt{Gm_p/R}$, 因而 $\Delta v=(1-\sqrt{3/2})\sqrt{Gm_p/R}<0$, 须减速. 初始远日速率还由角动量 $3Rv_2=Rv_p$ 确定, 不是独立自由参数.
\end{solution}
